\chapter*{Danh mục từ viết tắt và thuật ngữ}
\addcontentsline{toc}{chapter}{Danh mục từ viết tắt và thuật ngữ}
\markboth{DANH MỤC TỪ VIẾT TẮT}{}

\ghichu{Liệt kê các từ viết tắt và thuật ngữ xuất hiện trong báo cáo, sắp xếp theo bảng chữ cái. Bảng dưới đây đã điền sẵn các mục chắc chắn dùng đến; bổ sung thêm khi viết các chương sau.}
\nguon{docs/glossary.md}

\begin{longtable}{|p{3cm}|p{5cm}|p{6.5cm}|}
\hline
\textbf{Viết tắt} & \textbf{Thuật ngữ đầy đủ} & \textbf{Giải thích} \
\hline
\endfirsthead
\hline
\textbf{Viết tắt} & \textbf{Thuật ngữ đầy đủ} & \textbf{Giải thích} \
\hline
\endhead
API & Application Programming Interface & Giao diện lập trình ứng dụng \ \hline
CSDL & Cơ sở dữ liệu & Nơi lưu trữ dữ liệu có cấu trúc của hệ thống \ \hline
ERD & Entity Relationship Diagram & Sơ đồ quan hệ thực thể \ \hline
JWT & JSON Web Token & Chuẩn token dùng để xác thực người dùng \ \hline
RBAC & Role-Based Access Control & Kiểm soát truy cập theo vai trò \ \hline
REST & Representational State Transfer & Kiểu kiến trúc thiết kế dịch vụ web \ \hline
SPA & Single Page Application & Ứng dụng web một trang \ \hline
SSO & Single Sign-On & Đăng nhập một lần qua nhà cung cấp danh tính \ \hline
SVG & Scalable Vector Graphics & Định dạng ảnh vector dùng vẽ sơ đồ mặt bằng \ \hline
UAT & User Acceptance Testing & Kiểm thử chấp nhận người dùng \ \hline
\ldots & \ldots & \ldots \ \hline
\caption{Danh mục từ viết tắt sử dụng trong báo cáo}
\label{tab:viettat}
\end{longtable}
